\section*{Chapitre \ref{ch:b2:series} --- Suites et séries}

\begin{solution}{exo:b2:series:1}
$\sum\frac{\cos n}{n}$~: critère d'Abel avec $a_n = \frac1n$ et $b_n =
\cos n$, dont les sommes partielles sont bornées (partie réelle d'une
somme géométrique, comme dans l'\cref{ex:b2:series:sinn})~: convergente
(pas \omterm{def:b2:series:def}{absolument}, par la même astuce du $\cos^2$).

$\sum \frac{(-1)^n}{\ln n}$ ($n \geq 2$)~: critère des séries
\omterm{def:b2:linalg:alternating}{alternées}, $\frac{1}{\ln n} \downarrow 0$~: convergente~; pas
\omterm{def:b2:series:def}{absolument} ($\ln n \leq n$).

$\sum \frac{(-1)^n}{n^{3/4} + \cos n}$~: on développe,
\[
\frac{(-1)^n}{n^{3/4} + \cos n}
= \frac{(-1)^n}{n^{3/4}}\cdot
\frac{1}{1 + \frac{\cos n}{n^{3/4}}}
= \frac{(-1)^n}{n^{3/4}}
- \frac{(-1)^n\cos n}{n^{3/2}} + O\Bigl(\frac{1}{n^{9/4}}\Bigr).
\]
Première série~: \omterm{def:b2:linalg:alternating}{alternée}, convergente. Deuxième~: \omterm{def:b2:series:def}{absolument}
convergente (échelle $\frac{1}{n^{3/2}}$). Troisième~: \omterm{def:b2:series:def}{absolument}
convergente. Total~: convergente.
\end{solution}

\begin{solution}{exo:b2:series:2}
\omterm{thm:b2:series:fubini}{Produit de Cauchy} de $\sum_{m\geq0} z^m$ par elle-même (toutes deux
\omterm{def:b2:series:def}{absolument} convergentes pour $\abs z < 1$)~: le coefficient diagonal
est $c_k = \sum_{m=0}^{k} 1 = k + 1$, donc
\[
\frac{1}{(1-z)^2} = \sum_{k\geq0} (k+1)z^k .
\]
En multipliant par $z$ et en réindexant~: $\sum_{n \geq 1} n z^n =
\frac{z}{(1-z)^2}$.
\end{solution}

\begin{solution}{exo:b2:series:3}
Posons $B_n = \sum_{k \leq n} b_k \to B$. \omterm{thm:b2:series:abel}{Transformation d'Abel} avec
$a_k = k$~:
\[
\sum_{k=1}^{n} k\,b_k = n B_n - \sum_{k=1}^{n-1} B_k
\quad\Longrightarrow\quad
\frac1n \sum_{k=1}^{n} k b_k = B_n - \frac{1}{n}\sum_{k=1}^{n-1}
B_k .
\]
Les moyennes de Cesàro de la suite convergente $(B_k)$ tendent vers
sa limite $B$ (volume de Licence~1), donc le membre de droite tend
vers $B - B = 0$.
\end{solution}

\begin{solution}{exo:b2:series:4}
Si $\theta \in \pi\Z$~: tous les termes sont nuls --- convergente
trivialement. Supposons $\theta \notin \pi\Z$.

$\alpha > 1$~: \omterm{def:b2:series:def}{absolument} convergente (domination par $n^{-\alpha}$).

$0 < \alpha \leq 1$~: le critère d'Abel s'applique ($a_n = n^{-\alpha}
\downarrow 0$~; sommes partielles de $\sin n\theta$ bornées par
$\frac{1}{\abs{\sin(\theta/2)}}$, somme géométrique)~: convergente.
Pas \omterm{def:b2:series:def}{absolument}~: $\abs{\sin n\theta} \geq \sin^2 n\theta = \frac{1 -
\cos 2n\theta}{2}$, et $\sum \frac{1 - \cos 2n\theta}{2n^\alpha}$
diverge ($\sum n^{-\alpha}$ diverge~; $\sum
\frac{\cos 2n\theta}{n^\alpha}$ converge par Abel lorsque $2\theta
\notin 2\pi\Z$~; le cas exclu $2\theta \in 2\pi\Z$ signifie
$\theta \in \pi\Z$, déjà traité). Semi-convergente.
\end{solution}

\begin{solution}{exo:b2:series:5}
Soient $p_1, p_2, \dots$ les termes positifs de $(u_n)$ dans l'\omterm{def:b2:structures:generated}{ordre},
$q_1, q_2, \dots$ les termes négatifs. Les deux séries $\sum p_k$ et
$\sum q_k$ divergent~: si l'une convergeait, l'autre serait égale à
la série convergente $\sum u_n$ moins celle-ci, donc convergerait
aussi --- et alors $\sum \abs{u_n} = \sum p_k - \sum q_k$
convergerait, contredisant l'hypothèse. De plus $u_n \to 0$ ($\sum
u_n$ converge).

Réordonnancement glouton~: prendre des termes positifs $p_1, p_2,
\dots$ jusqu'à ce que le total courant dépasse $\ell$ pour la
première fois (possible~: $\sum p_k = +\infty$)~; puis des termes
négatifs jusqu'à ce que le total passe sous $\ell$ pour la première
fois (possible~: $\sum q_k = -\infty$)~; répéter indéfiniment (chaque
phase est finie, et chaque terme est utilisé exactement une fois~: un
vrai réordonnancement). Après chaque changement, la distance du total
courant à $\ell$ vaut au plus le dernier terme utilisé~; comme les
termes utilisés au $m$-ième changement ont un indice $\to \infty$, et
que $u_n \to 0$, les totaux courants convergent vers $\ell$.
\end{solution}

\begin{solution}{exo:b2:series:6}
Sommabilité~: les sommes partielles finies de
$\frac{\abs x^{m+n}}{m!n!}$ sont bornées par $\bigl(\sum_m
\frac{\abs x^m}{m!}\bigr)^2 = \eu^{2\abs x}$. Regroupement diagonal
(\cref{thm:b2:series:fubini})~:
\[
(\eu^{x})^2 = \sum_{m,n} \frac{x^{m+n}}{m!\,n!}
= \sum_{k=0}^{\infty} x^k \sum_{m+n=k} \frac{1}{m!\,n!}
= \sum_k \frac{x^k}{k!}\sum_{m=0}^{k}\binom km
= \sum_k \frac{(2x)^k}{k!} = \eu^{2x} .
\]
\end{solution}

\begin{solution}{exo:b2:series:7}
Sommabilité~: bornée par $\zeta(2)^2$ en tant que famille produit
(l'argument du \omterm{thm:b2:series:fubini}{produit de Cauchy} du~\cref{thm:b2:series:fubini}).
L'application $(q, a, b) \mapsto (qa, qb)$, des triplets avec
$\gcd(a, b) = 1$ vers les couples $(m, n)$, est une bijection (poser
$q = \gcd(m,n)$). En regroupant la \omterm{def:b2:series:summable}{famille sommable} en conséquence
(une partition de l'ensemble d'indices --- légitime pour les \omterm{def:b2:series:summable}{familles
sommables} d'après le~\cref{thm:b2:series:rearrangement}/le~\cref{thm:b2:series:fubini}
appliqué à la partition en une infinité \omterm{def:b2:structures:countable}{dénombrable} de classes)~:
\[
\zeta(2)^2 = \sum_{q} \sum_{\gcd(a,b)=1} \frac{1}{q^4 a^2 b^2}
= \zeta(4)\, S,
\qquad\text{donc}\qquad
S = \frac{\zeta(2)^2}{\zeta(4)} .
\]
(Avec les valeurs $\zeta(2) = \frac{\pi^2}{6}$,
$\zeta(4) = \frac{\pi^4}{90}$ du \cref{ch:b2:fourier}~: $S =
\frac{5}{2}$.)
\end{solution}

\begin{solution}{exo:b2:series:8}
Posons $A = \sum a_n$ (absolue), $B_m = \sum_{k\leq m} b_k \to B$,
bornée par $M$. Alors
\[
C_N = \sum_{k=0}^{N} c_k = \sum_{n=0}^{N} a_n B_{N-n}
\]
(en regroupant par l'indice de $a$). Écrivons
\[
C_N - AB = \sum_{n=0}^{N} a_n (B_{N-n} - B) - B\sum_{n > N} a_n .
\]
Le dernier terme tend vers $0$. Scindons la somme en $n = \lfloor N/2
\rfloor$~: pour $n \leq N/2$, $N - n \geq N/2$, donc $\abs{B_{N-n} - B}
\leq \varepsilon_N := \sup_{m \geq N/2}\abs{B_m - B} \to 0$, et
cette partie est $\leq \varepsilon_N \sum\abs{a_n}$~; pour $n > N/2$,
$\abs{B_{N-n} - B} \leq 2M$, et cette partie est $\leq 2M \sum_{n >
N/2} \abs{a_n} \to 0$. D'où $C_N \to AB$.
\end{solution}

\begin{solution}{exo:b2:series:9}
Prenons $u_n = 2^{-n} e_n$ ($e_n$ la suite ayant un unique $1$ en
position $n$). Alors $\sum \norm{u_n}_\infty = \sum 2^{-n} <
\infty$~: \omterm{def:b2:series:def}{absolument} convergente. Mais les sommes partielles $S_N =
(1, \tfrac12, \dots, 2^{-N}, 0, \dots)$ devraient converger vers la
suite $(2^{-n})_n$, qui n'est \emph{pas} nulle à partir d'un certain
rang~: en dehors de $E$. Dans $E$, $(S_N)$ est de Cauchy sans limite
($\norm{S_N - x}_\infty \geq 2^{-N-1}$ ne peut aider aucun $x$ nul à
partir d'un certain rang~: pour tout $x \in E$ s'annulant au-delà du
rang $K$, $\norm{S_N - x} \geq 2^{-K-1}$ pour $N > K$)~: la série ne
converge pas dans $E$. La complétude est exactement ce dont
le~\cref{thm:b2:nvs:absoluteconvergence} a besoin.
\end{solution}

\begin{solution}{exo:b2:series:10}
$\arctan(n+1) - \arctan n$ et $\arctan\frac{1}{n^2+n+1}$ appartiennent
tous deux à $\intoo{0}{\frac\pi2}$, et la formule d'addition de la
tangente donne
\[
\tan\bigl(\arctan(n{+}1) - \arctan n\bigr)
= \frac{(n+1) - n}{1 + n(n+1)} = \frac{1}{n^2 + n + 1} :
\]
tangentes égales dans un intervalle où $\tan$ est injective, donc
l'identité est vraie. Par télescopage,
\[
\sum_{n=1}^{N}\arctan\frac{1}{n^2+n+1}
= \arctan(N{+}1) - \arctan 1 \xrightarrow[N\to\infty]{}
\frac\pi2 - \frac\pi4 = \frac\pi4 .
\]

\end{solution}

\begin{solution}{exo:b2:series:11}
Première~: $\sqrt{n+1} - \sqrt n = \frac{1}{\sqrt{n+1} + \sqrt n}
\sim \frac{1}{2\sqrt n}$, donc les termes sont $\sim
2^{-\alpha}n^{-\alpha/2}$~: convergence si et seulement si
$\frac\alpha2 > 1$, c'est-à-dire $\alpha > 2$. Deuxième~: $1 -
\cos\frac1n \sim \frac{1}{2n^2}$~: converge. Troisième~: $\bigl(1 +
\frac1n\bigr)^n = \eu^{\,n\ln(1 + 1/n)} = \eu^{\,1 - \frac1{2n} +
O(n^{-2})} = \eu\bigl(1 - \frac{1}{2n} + O(n^{-2})\bigr)$, donc
\[
\eu - \Bigl(1 + \frac1n\Bigr)^{\!n} \sim \frac{\eu}{2n} :
\]
termes positifs équivalents à un multiple harmonique~: diverge.
\end{solution}

\begin{solution}{exo:b2:series:12}
Par monotonie, $n\,a_{2n} \leq a_{n+1} + a_{n+2} + \dots +
a_{2n} = S_{2n} - S_n \to 0$ (critère de Cauchy pour la série
convergente). D'où $2n\,a_{2n} \to 0$, et $(2n{+}1)\,
a_{2n+1} \leq (2n{+}1)a_{2n} = \frac{2n+1}{2n}\,(2n\,a_{2n}) \to
0$~: les deux sous-suites de $(na_n)$ tendent vers $0$, donc $na_n
\to 0$.

\emph{La réciproque est fausse~:} $a_n = \frac1{n\ln n}$ est positive
décroissante avec $na_n = \frac1{\ln n} \to 0$, et pourtant $\sum a_n$
diverge (frontière de Bertrand, \cref{pb:b2:series:1}, question~18).
\emph{Monotonie essentielle~:} posons $a_n = \frac1n$ quand $n$ est
une puissance de $2$ et $a_n = 2^{-n}$ sinon~: $\sum a_n \leq
\sum_k 2^{-k} + \sum_n 2^{-n} < \infty$, mais $na_n = 1$ le long des
puissances de $2$~: $na_n \not\to 0$.
\end{solution}

\begin{solution}{pb:b2:series:1}
\textbf{1.} Récurrence sur $m$~: pour $m = 0$ les deux membres valent
$1$~; l'étape multiplie par $\cos\theta + \iu\sin\theta$ et utilise
les formules d'addition $\cos(m\theta + \theta) =
\cos m\theta\cos\theta - \sin m\theta\sin\theta$, $\sin(m\theta
+ \theta) = \sin m\theta\cos\theta + \cos m\theta\sin\theta$. En
développant plutôt par la formule du binôme avec $m = 2n+1$ et en
collectant la partie imaginaire (les puissances impaires de
$\iu\sin\theta$, avec $\iu^{2j+1} = (-1)^j\iu$)~:
\[
\sin\bigl((2n{+}1)\theta\bigr) = \sum_{j=0}^{n}(-1)^j
\binom{2n+1}{2j+1}\cos^{2(n-j)}\theta\,\sin^{2j+1}\theta .
\]

\textbf{2.} Sur $\intoo0{\frac\pi2}$, $\sin\theta \neq 0$~: on
factorise $\sin^{2n+1}\theta$ dans chaque terme, laissant
$\bigl(\frac{\cos^2\theta}{\sin^2\theta}\bigr)^{n-j} =
(\cot^2\theta)^{n-j}$~: l'identité affichée avec $P_n(x) =
\sum_j(-1)^j\binom{2n+1}{2j+1}x^{n-j}$. Son coefficient de degré $n$
est celui de $j = 0$, à savoir $\binom{2n+1}{1} = 2n + 1 \neq 0$.

\textbf{3.} En $\theta_k = \frac{k\pi}{2n+1}$~:
$\sin\bigl((2n{+}1)\theta_k\bigr) = \sin k\pi = 0$ alors que
$\sin^{2n+1}\theta_k \neq 0$, donc $P_n(\cot^2\theta_k) = 0$. Les
$\theta_k$ croissent strictement dans $\intoo0{\frac\pi2}$, où
$\cot^2$ est strictement décroissante~: les valeurs $x_k =
\cot^2\theta_k$ sont deux à deux distinctes --- $n$ racines distinctes
d'un polynôme de degré $n$, donc toutes ses racines.

\textbf{4.} Viète~: la somme des racines est l'opposé du rapport des
coefficients de $x^{n-1}$ et $x^n$~:
\[
\sum_{k=1}^{n}\cot^2\theta_k =
\frac{\binom{2n+1}{3}}{\binom{2n+1}{1}}
= \frac{(2n+1)(2n)(2n-1)/6}{2n+1} = \frac{n(2n-1)}{3}.
\]

\textbf{5.} $\frac{1}{\sin^2\theta} = 1 + \cot^2\theta$~: en sommant,
$n + \frac{n(2n-1)}3 = \frac{3n + 2n^2 - n}{3} =
\frac{2n(n+1)}{3}$.

\textbf{6.} Sur $\intoo{0}{\frac\pi2}$~: $\sin\theta < \theta <
\tan\theta$ (volume de Licence~1). Le passage aux inverses renverse~:
$\cot\theta < \frac1\theta < \frac1{\sin\theta}$, et l'élévation au
carré (tout est positif) donne $\cot^2\theta < \frac1{\theta^2} <
\frac1{\sin^2\theta}$.

\textbf{7.} Sommons la question~6 en $\theta = \theta_k$ sur $k \leq
n$, en utilisant les questions~4 et 5, et $\frac1{\theta_k^2} =
\frac{(2n+1)^2}{k^2\pi^2}$~:
\[
\frac{n(2n-1)}3 < \frac{(2n+1)^2}{\pi^2}\sum_{k=1}^n\frac1{k^2}
< \frac{2n(n+1)}3 .
\]

\textbf{8.} Multiplions par $\frac{\pi^2}{(2n+1)^2}$~:
\[
\frac{\pi^2}{3}\cdot\frac{n(2n-1)}{(2n+1)^2}
< \sum_{k=1}^{n}\frac1{k^2} <
\frac{\pi^2}{3}\cdot\frac{2n(n+1)}{(2n+1)^2}.
\]
Les deux bornes tendent vers $\frac{\pi^2}3\cdot\frac12 =
\frac{\pi^2}6$ (les fractions rationnelles tendent vers $\frac12$).
Les sommes partielles croissent, donc elles convergent, et
l'encadrement donne $\zeta(2) = \frac{\pi^2}6$~: le théorème d'Euler,
par la preuve de Cauchy.

\textbf{9.} Les sommes partielles croissent vers $\zeta(2) =
\frac{\pi^2}6$, donc $0 \leq \frac{\pi^2}6 - \sum_{k\leq n}k^{-2}$~;
et la borne inférieure de la question~8 donne
\[
\frac{\pi^2}6 - \sum_{k\leq n}\frac1{k^2}
\leq \frac{\pi^2}6 - \frac{\pi^2}3\cdot\frac{n(2n-1)}{(2n+1)^2}
= \frac{\pi^2}6\cdot\frac{(2n+1)^2 - (4n^2 -
2n)}{(2n+1)^2}
= \frac{\pi^2}6\cdot\frac{6n + 1}{(2n+1)^2}
= O\Bigl(\frac1n\Bigr),
\]
en accord avec le reste exact $\frac1n - \frac1{2n^2} + O(n^{-3})$ de
\cref{exo:b2:comparison:11}.

\textbf{10.} La deuxième fonction symétrique élémentaire des racines
est $\sigma_2 = \frac{\binom{2n+1}5}{\binom{2n+1}1} =
\frac{(2n)(2n-1)(2n-2)(2n-3)}{120}$, donc
\[
\sum_k\cot^4\theta_k = \sigma_1^2 - 2\sigma_2
= \Bigl(\frac{n(2n-1)}3\Bigr)^2 -
\frac{2n(2n-1)(2n-2)(2n-3)}{60}
\sim \frac{4n^4}9 - \frac{4n^4}{15} = \frac{8n^4}{45}.
\]
En encadrant $\cot^4\theta < \theta^{-4} < (1 + \cot^2\theta)^2 =
1 + 2\cot^2\theta + \cot^4\theta$ et en sommant~: les deux sommes
extérieures valent $\frac{8n^4}{45}(1 + o(1))$ (le terme ajouté $n +
2\sigma_1 = O(n^2)$ est négligeable), tandis que celle du milieu est
$\frac{(2n+1)^4}{\pi^4}\sum_{k\leq n}k^{-4}$. D'où
\[
\sum_{k\leq n}\frac1{k^4} \longrightarrow
\pi^4\cdot\frac{8/45}{16} = \frac{\pi^4}{90}.
\]

\textbf{11.} En scindant $\zeta(2)$ selon les parités~:
$\sum_{\text{pair}} = \sum_j\frac{1}{(2j)^2} = \frac14\zeta(2) =
\frac{\pi^2}{24}$, donc $\sum_{\text{impair}} = \zeta(2) -
\frac{\pi^2}{24} = \frac{\pi^2}8$. \omterm{def:b2:linalg:alternating}{Alternée}~:
$\sum_k\frac{(-1)^{k-1}}{k^2} = \sum_{\text{impair}} -
\sum_{\text{pair}} = \frac{\pi^2}8 - \frac{\pi^2}{24} =
\frac{\pi^2}{12}$ (la convergence absolue justifie le regroupement,
\cref{thm:b2:series:rearrangement}).

\textbf{12.} $S = \frac{\zeta(2)^2}{\zeta(4)} =
\frac{(\pi^2/6)^2}{\pi^4/90} = \frac{90}{36} = \frac52$.
Heuristique~: l'identité $\zeta(2)^2 = \zeta(4)S$ de
l'\cref{exo:b2:series:7} dit qu'extraire le $\gcd$ renormalise les
couples en couples premiers entre eux~; l'inverse
$\frac1{\zeta(2)} = \frac6{\pi^2} \approx 0.608$ est le candidat
naturel pour la densité des couples premiers entre eux parmi tous les
couples --- un énoncé sur $\lim_N \frac{1}{N^2}\#\{(m,n) \leq N :
\gcd = 1\}$ dont la preuve honnête (avec termes d'erreur) relève du
volume de Licence~3.

\textbf{13.} La sommation directe a une erreur $\sim \frac1n$~: six
chiffres exigent environ $10^6$ termes. La somme corrigée
$\sum_{k\leq n}k^{-2} + \frac1n - \frac1{2n^2}$ a une erreur
$O(n^{-3})$~: en $n = 100$ elle vaut $1.6449339\dots$ contre
$\frac{\pi^2}6 = 1.6449341\dots$ --- erreur $1.7\cdot10^{-7}$,
sept chiffres à partir de cent termes. Les corrections asymptotiques
battent la patience brute de quatre \omterm{def:b2:structures:generated}{ordres} de grandeur.

\textbf{14.} $n = 1$~: $P_1(x) = 3x - 1$, racine $\frac13$, et
en effet $\cot^2\frac\pi3 = \bigl(\frac1{\sqrt3}\bigr)^2 =
\frac13 = \frac{1\cdot1}3$. $n = 2$~: la formule prédit
$\frac{2\cdot3}3 = 2$~; avec $\cos\frac\pi5 = \frac{1 +
\sqrt5}{4}$, on calcule $\cot^2 36^\circ \approx 1.894$ et
$\cot^2 72^\circ \approx 0.106$~: somme $2.000$.

\textbf{15.} Les nombres $\tan^2\theta_k = \frac1{x_k}$ sont les
racines de $Q(x) = x^nP_n\bigl(\frac1x\bigr) =
\sum_{j=0}^n(-1)^j\binom{2n+1}{2j+1}x^j$ (les $x_k$ sont non nuls).
Viète sur $Q$~: le coefficient dominant est $(-1)^n$ (terme $j = n$),
le suivant est $(-1)^{n-1}\binom{2n+1}{2n-1} =
(-1)^{n-1}\binom{2n+1}{2}$, donc
\[
\sum_{k=1}^n\tan^2\theta_k =
-\frac{(-1)^{n-1}\binom{2n+1}2}{(-1)^n} = \binom{2n+1}2\cdot
\frac{2}{2n+1}\cdot\frac{2n+1}{2}
= n(2n+1).
\]
Vérification $n = 1$~: $\tan^2\frac\pi3 = 3 = 1\cdot3$.

\textbf{16.} Posons $\gamma = \frac{1+\alpha}2 \in
\intoo{1}{\alpha}$. Alors $\frac{n^{-\alpha}(\ln
n)^{-\beta}}{n^{-\gamma}} = n^{\gamma - \alpha}(\ln n)^{-\beta}
\to 0$ (une puissance négative de $n$ l'emporte sur toute puissance
de $\ln n$), donc à partir d'un certain rang les termes sont $\leq
n^{-\gamma}$ avec $\gamma > 1$~: convergence par comparaison avec une
série de Riemann.

\textbf{17.} Posons $\gamma = \frac{1+\alpha}2 \in
\intoo{\alpha}{1}$~: cette fois $\frac{n^{-\gamma}}{n^{-\alpha}(\ln
n)^{-\beta}} = n^{\alpha-\gamma}(\ln n)^{\beta} \to 0$, donc à partir
d'un certain rang les termes sont $\geq n^{-\gamma}$ avec $\gamma <
1$~: divergence.

\textbf{18.} $f(t) = \frac1{t(\ln t)^\beta}$ est positive, \omterm{def:b2:metric:continuity}{continue},
et décroissante pour $t$ grand (son logarithme a pour dérivée
$-\frac1t\bigl(1 + \frac{\beta}{\ln t}\bigr) < 0$ à partir d'un
certain rang). Primitives~: pour $\beta \neq 1$,
$\int^x f = \frac{(\ln x)^{1-\beta}}{1-\beta} + \text{cste}$, qui a
une limite finie si et seulement si $\beta > 1$~; pour $\beta = 1$,
$\int^x f = \ln\ln x \to \infty$. D'après
le~\cref{thm:b2:comparison:seriesintegral}, la série et l'intégrale
ont la même nature~: convergence si et seulement si $\beta > 1$.

\textbf{19.} Même critère~: $\frac{\dd}{\dd t}\ln\ln\ln t =
\frac{1}{t\ln t\,\ln\ln t}$, et $\ln\ln\ln t \to \infty$~:
divergence. Et $\frac{\dd}{\dd t}\Bigl(-\frac1{\ln\ln t}\Bigr)
= \frac{1}{t\ln t\,(\ln\ln t)^2}$ avec $-\frac1{\ln\ln t} \to
0$~: convergence.

\textbf{20.} Première~: $n^{1/\ln n} = \eu^{\ln n/\ln n} = \eu$, donc
les termes valent exactement $\frac{1}{\eu\,n}$~: un multiple de la
série harmonique, \emph{divergente} --- l'exposant $1 +
\frac1{\ln n}$ rampe vers $1$ trop vite. Deuxième~: $n^{1/\ln\ln n}
= \eu^{\ln n/\ln\ln n}$, et $\frac{\ln n}{\ln\ln n} \geq
2\ln\ln n$ à partir d'un certain rang, donc $n^{1/\ln\ln n} \geq
(\ln n)^2$~: les termes sont $\leq \frac1{n(\ln n)^2}$, une série de
Bertrand convergente (question~18)~: \emph{convergente}. La frontière
passe strictement entre ces deux exposants.

\textbf{21.} $R_n \downarrow 0$ et
\[
\sqrt{R_n} - \sqrt{R_{n+1}} = \frac{R_n -
R_{n+1}}{\sqrt{R_n} + \sqrt{R_{n+1}}} =
\frac{a_n}{\sqrt{R_n} + \sqrt{R_{n+1}}} \geq
\frac{a_n}{2\sqrt{R_n}},
\]
donc $\sum_n \frac{a_n}{\sqrt{R_n}} \leq 2\sum_n(\sqrt{R_n} -
\sqrt{R_{n+1}}) = 2\sqrt{R_1} < \infty$ (télescopage). Et pourtant
$\frac{a_n/\sqrt{R_n}}{a_n} = \frac1{\sqrt{R_n}} \to \infty$~: la
nouvelle série converge tout en étant infiniment plus grande. Aucune
série convergente n'est la plus lente~; les tests de comparaison
contre une famille fixée quelconque ne peuvent jamais être \omterm{def:b2:metric:complete}{complets}.

\textbf{22.} Pour $(a_n)$ positive décroissante, regroupons les
termes entre puissances consécutives de $2$~:
\[
2^{k}a_{2^{k+1}} \leq \sum_{n=2^k}^{2^{k+1}-1} a_n \leq
2^ka_{2^k} .
\]
En sommant sur $k$~: si $\sum 2^ka_{2^k}$ converge, les sommes
partielles de $\sum a_n$ sont bornées (converge)~; si $\sum a_n$
converge, alors $\sum_k 2^{k+1}a_{2^{k+1}} \leq 2\sum_n a_n <
\infty$. Pour $a_n = \frac1{n(\ln n)^\beta}$~: $2^ka_{2^k} =
\frac{1}{(k\ln 2)^\beta}$, et $\sum k^{-\beta}$ converge si et
seulement si $\beta > 1$~: de nouveau la frontière de la question~18,
sans intégrales.

\textbf{23.} Par la comparaison intégrale,
\[
\sum_{n > N}\frac{1}{n(\ln n)^2} \geq
\int_{N+1}^{\infty}\frac{\dd t}{t(\ln t)^2} =
\frac{1}{\ln(N+1)},
\]
qui en $N = 10^6$ vaut $\approx 0.0724$~: après un million de termes
le reste dépasse encore $0.07$ --- la série converge, mais aucune
sommation directe n'en exhibera jamais la somme. À contraster avec la
question~13, où une seule correction asymptotique a acheté sept
chiffres à partir de cent termes~: savoir \emph{comment} une série
converge vaut mieux que savoir qu'elle converge.

\textbf{24.} $\sum\frac1{n\ln n}$~: diverge ($\alpha = 1$,
$\beta = 1$, question~18). $\sum\frac1{n^{1.01}}$~: converge
(Riemann, $\alpha > 1$). $\sum\frac{(\ln n)^{100}}{n^{1.001}}$~:
converge ($\alpha = 1.001 > 1$, $\beta = -100$, question~16).
$\sum\frac1{n\ln n(\ln\ln n)^3}$~: converge (schéma de la
question~19~: primitive $-\frac12(\ln\ln t)^{-2}$, limite finie).

\textbf{25.} De Moivre convertit l'annulation de $\sin(2n{+}1)
\theta_k$ en l'annulation d'un polynôme en $\cot^2\theta_k$, et Viète
y lit les sommes de racines exactes que l'analyse seule ne pouvait
qu'estimer (questions~1--5). L'encadrement a eu besoin des valeurs
\emph{exactes} $\frac{n(2n-1)}3$ et $\frac{2n(n+1)}3$ des deux côtés
--- des équivalents auraient été une pétition de principe, puisque
tout l'enjeu est la constante $\frac{\pi^2}6$ (questions~7--8). La
partie~IV a entièrement reposé sur la comparaison série--intégrale du
\cref{ch:b2:comparison}, les primitives logarithmiques faisant le
classement (questions~18--19). La question~21 détruit le rêve d'un
test de comparaison universel~: sous toute série convergente s'en
trouve une autre, infiniment plus lente --- des échelles comme celle
de Bertrand cartographient la frontière toujours plus finement mais
ne l'atteignent jamais. Sommets~: le $\zeta(2)
= \frac{\pi^2}6$ d'Euler avec son étage supérieur $\zeta(4) =
\frac{\pi^4}{90}$ (questions~8, 10), et la classification de Bertrand
(questions~16--18)~; $\zeta(2)$ revient au \cref{ch:b2:fourier}, où
l'identité de Parseval le reprouve en une ligne à partir de la série
de Fourier du signal en dents de scie --- une seule constante, deux
civilisations.
\end{solution}
